\textcolor{red}{(Este capítulo em geral não deve ultrapassar o total de 5 páginas).}

\section{Considerações Iniciais}

\textcolor{red}{Neste capítulo devem ser apresentados os conceitos e a terminologia básicos da área do projeto e o levantamento bibliográfico necessário para a realização do trabalho. Em particular, a descrição dos principais trabalhos de pesquisa relacionados com este (em geral, aqueles que representam o estado da arte na área de pesquisa em questão), bem como dos trabalhos que serviram de base para a solução proposta por este projeto (em geral, aqueles que apresentam técnicas ou recursos que foram utilizados pelo projeto). A divisão nas subseções a seguir é opcional. Em particular, nesta seção (Considerações Iniciais), deve-se descreva sucintamente o que é apresentado neste capítulo.}

\section{Título}
\subsection{Subtítulo 1}
\subsection{Subtítulo 2}
\textcolor{red}{(As seções e subsubseções podem ser organizadas da forma que melhor se adequar à sua monografia)}

\section{Trabalhos Relacionados}

\textcolor{red}{Refereciar aqui os trabalhos relacionados ao projeto. Não esquecer de citar a fonte corretamente, por exemplo [Silva e Zhao 2012].}

\section{Considerações Finais}

\textcolor{red}{Nesta seção deve-se apresentar uma conclusão sobre este capítulo e introduzir brevemente o capítulo seguinte.}